\section{背景与目标}

在医疗、金融等高敏感场景中，数据使用方通常不愿以明文形式交出原始数据，模型提供方也不愿直接暴露明文模型参数。单向隐私保护机制难以同时满足双方诉求：只强调输入侧保护会暴露模型知识产权，只强调模型侧保护则会引发数据合规风险。进一步地，若模型采用按样本动态变化的剪枝路径，则依赖明文删除、阈值比较和变长前向结构的决策链很难直接迁入安全计算图，这也是许多安全推理系统被迫退化为固定结构模型的重要原因。

本作品的设计目标不是单纯提升胸片分类精度，而是在双向隐私前提下保留动态剪枝能力，并将其落实到可运行的工程系统中。为此，本作品选择动态词元剪枝视觉 Transformer 方法 DynamicViT 作为主要载体，围绕词元打分、阈值比较、保留决策与安全执行语义之间的映射关系，完成从模型结构、协议路径到前后端控制面的联动设计。

从问题形式看，本作品研究的是一个典型的两方私有推理任务：数据持有方希望在不泄露原始影像的前提下获得分类结果，模型持有方希望在不公开模型参数与中间特征的前提下完成服务。若将原始模型记为 f、输入影像记为 x，系统需要在两方仅持有各自秘密份额的条件下近似完成 f(x) 的计算，并限制任何单方从通信内容、中间状态或动态路径中恢复对方敏感信息。该形式化目标决定了作品不能只把明文模型封装为 Web 服务，而必须重新审视动态剪枝、控制流、张量形状和审计元数据在安全计算中的暴露边界。

\section{应用场景与交付边界}

本作品面向医院侧数据与模型提供方参数同时敏感的协同推理场景，目标是在不暴露原始影像和明文模型参数的前提下完成医疗影像分类。系统的主要交付对象包括终端用户、医院侧部署环境、服务端控制面以及两方安全执行域。终端侧负责本地解码、预处理、质量评估与秘密分享生成；服务端负责协议预检、重放与并发守卫、审计落盘和密态推理调度；安全执行域负责完成双向隐私约束下的前向计算并返回最终分类结果。

本作品的交付边界不包括生产级拒绝服务防护、恶意参与方串谋对抗、输出侧隐私泄露建模以及完整的多中心部署评估。本文的重点是证明：在双向隐私约束下，动态剪枝语义能够被安全化重写，并形成可运行、可验证、可复现实验闭环。

从隐私保护推理研究脉络看，早期系统多围绕卷积网络或一般前馈模型展开，例如 SecureML[14]、GAZELLE[15]、CrypTen[16] 和 Delphi[17] 等工作已经证明，在半诚实模型下可以同时保护输入与模型参数，并将主要优化集中在线性层、激活函数和加密切换流程上。然而，当模型迁移到 Transformer 后，注意力、层归一化（Layer Normalization, LayerNorm）、高斯误差线性单元（Gaussian Error Linear Unit, GELU）和大规模矩阵乘使通信量、交互轮次与数值稳定性问题同时放大，原有方案难以直接复用到动态词元剪枝场景中。

对医疗影像任务而言，问题又比通用私有推理更复杂。一方面，医院侧原始影像具有强隐私属性，模型提供方也往往不愿公开完整参数；另一方面，胸片等视觉任务存在明显的 patch 冗余，如果为了协议简化完全退回固定结构模型，又会削弱动态剪枝在时延和通信量上的潜在优势。因此，本作品关注的不是一般意义上的“能否安全推理”，而是在双向隐私、动态边界保留与工程闭环交付三者之间建立可验证的平衡。

\section{核心贡献}

本作品的核心贡献主要体现在三个方面。第一，在算法层面，将 DynamicViT 中依赖明文控制流的删除式剪枝路径重写为固定图下可执行的密态语义，使动态词元剪枝能够进入安全计算框架。第二，在系统层面，将浏览器工作线程、服务端权威快检、重放与并发守卫以及审计链整合为完整控制面闭环，保证输入、传输和执行三个阶段的边界清晰可验证。第三，在交付层面，给出可复现的实验流程、基线对比与消融验证，使系统不仅“能跑”，而且能够说明“为什么这样设计是必要的”。

\section{相关工作与技术定位}

在隐私保护视觉 Transformer 推理方向，已有研究开始关注如何在不暴露用户输入的前提下完成密态分类任务。其中，基于秘密共享的轻量级隐私保护 ViT 推理框架 SViT 是一条具有代表性的技术路线。该工作采用两个边缘服务器协同执行安全计算协议，并围绕 ViT-B/16 的关键非线性模块构建了 SExp、SDiv、SSqrt、SLayerNorm、SSoftmax、SGeLU 等协议，在保持原始 ViT-B/16 结构不变的前提下完成安全推理[13]。对于 ViT 路线，其核心注意力计算通常可写为式（1-1）：

该类工作的主要启发在于：对于隐私保护 ViT 系统，性能优化不一定只能依赖模型压缩或结构删减，也可以通过专门设计安全算子协议来降低在线计算与通信开销。文献[13]给出的结果表明，相比 CrypTen[16]，其框架在计算效率和在线通信开销方面分别可提升约 2~6 倍和 4~14 倍，这说明“围绕算子协议进行优化”本身就是一条有效路线。

但与本作品相比，SViT 的目标和系统边界并不完全相同。SViT 更强调在固定 ViT 结构下优化安全协议效率，采用的是两个边缘服务器加可信第三方离线辅助的体系；而本作品关注的是在不引入额外可信第三方的前提下，实现面向 DynamicViT 的双向隐私推理，并将动态词元剪枝边界重写为固定图下可执行的密态语义，同时补齐浏览器工作线程、服务端权威快检与审计链闭环。因此，SViT 更适合作为本作品的相关工作与技术参照，而不是对本作品主线的直接替代。

从更长的研究脉络看，隐私保护推理最早并非围绕 Transformer 展开，而是先在卷积网络和一般前馈网络上建立问题定义与协议基础。MiniONN[25]强调将已经训练完成的神经网络转换为 oblivious neural network，使既有模型在不改训练流程的前提下进入私有推理；XONN[26]则进一步利用二值化网络与 garbled circuits 显著降低线性层代价，展示了从“直接保护现有模型”到“为安全推理重塑模型形态”的演进趋势。结合 SecureML[14]、GAZELLE[15]、Delphi[17] 与 CrypTen[16] 可以看到，早期工作已经较好解决了输入保护、线性层计算和通用推理接口等基础问题，但它们大多默认网络结构固定、中间控制流简单，尚未把样本相关动态路径视为需要重点保护的对象。

当研究对象从卷积网络转向 Transformer 后，困难的焦点也随之转移。相较卷积结构，Transformer 不仅具有更高维度的矩阵乘、更频繁的归一化与激活计算，而且其实际开销会随 token 数量显著放大。因此，Transformer 私有推理的核心不再只是“如何安全地实现已有算子”，还包括“如何让模型结构本身更适合密态执行”。这也是后续 Iron[18]、BOLT[19]、BumbleBee[7]、NEXUS[20]、THE-X[22]、MPCFormer[23] 与 SecFormer[24] 等工作持续出现的重要背景。

除 SViT[13] 外，私有推理系统还经历了从通用神经网络到 Transformer 专用协议的演进。SecureML[14]、MiniONN[25]、GAZELLE[15]、XONN[26]、CrypTen[16] 与 Delphi[17] 分别从秘密共享、混合协议、garbled circuits、系统接口与工程抽象等方面推动了通用私有推理的发展；进入 Transformer 时代后，Iron[18]首先系统性讨论了 Softmax、GELU、LayerNorm 和大矩阵乘在私有推理中的复合开销，说明 Transformer 专用协议优化是独立问题，而不是卷积网络方案的简单平移。

在固定结构安全 Transformer 推理方向，现有研究大致可以分为三类。第一类以 BOLT[19] 为代表，强调在交互式协议环境中优化矩阵乘与非线性算子，从而在保留模型精度的前提下降低通信与运行成本。第二类以 BumbleBee[7] 与 THE-X[22] 为代表，更强调对大模型或同态近似路径的系统整合，通过协议组合、算子替换和近似函数设计压缩整体执行代价。第三类则以 NEXUS[20] 和 SecFormer[24] 为代表，进一步关注交互轮次与算子形态对在线性能的影响，尝试通过非交互式执行或 SMPC 友好近似来降低实际部署成本。上述研究虽然路线不同，但共同前提是模型结构大体固定，优化重点仍然落在算子与协议层。

另一类与本作品关系更直接的研究来自明文高效视觉 Transformer 领域。DynamicViT[2] 通过逐阶段词元打分和删除式稀疏化降低后续层计算量；EViT[27] 通过保留高关注 token 并融合低关注 token 来减少冗余计算；A-ViT[28] 借助自适应停止机制使不同 token 在不同深度提前退出；ToMe[29] 则通过 token merging 在几乎不改训练流程的前提下提高吞吐；MPCViT[3] 进一步把“模型结构是否适合 MPC 执行”显式纳入搜索目标。这些方法说明，视觉 Transformer 的效率问题与 token 数量管理高度相关，但它们默认剪枝、融合、停止或合并决策可以在明文环境下直接观察和执行。

正因为如此，明文高效 ViT 方法通常不能无改动地迁移到私有推理场景。无论是删除 token、让 token 提前停止，还是合并相似 token，本质上都涉及数据相关路径选择、变长结构或显式 importance 排序；若直接在安全执行域外部完成这些操作，就可能泄露输入分布、模型偏好甚至部分中间状态；若完全保留原始明文逻辑，又会使密态图难以固定、通信量和比较开销迅速上升。CipherPrune[21] 开始把词元重要性差异直接纳入私有推理优化目标，通过密态渐进式剔除低重要性词元并配合不同强度的多项式近似，同时缓解 token 数量和非线性算子开销带来的双重成本。与这些工作相比，本作品并不试图重建一套新的通用 Transformer 近似体系，而是选择以医疗影像 DynamicViT 为具体载体，把删除式词元剪枝安全化重写为 keep-mask 驱动的固定图语义，并在两方 SPU 执行、浏览器本地分片、服务端权威快检与审计闭环之间保持一致的工程边界。

综合来看，现有相关工作大致可分为三条路线：一是以 SecureML、MiniONN、GAZELLE、XONN、Delphi 和 CrypTen 为代表的通用私有推理系统，重点解决基础协议与通用神经网络算子问题；二是以 Iron、BOLT、BumbleBee、NEXUS、THE-X、MPCFormer、SecFormer 为代表的固定结构 Transformer 私有推理路线，重点降低 attention 及非线性算子的密态执行开销；三是以 DynamicViT、EViT、A-ViT、ToMe 和 CipherPrune 为代表的 token 稀疏化或动态推理路线，重点减少后续层 token 规模。本作品位于后两条路线的交叉点：它既继承了视觉 Transformer 动态 token 管理的建模思想，又必须满足两方私有推理的固定图与边界保护约束，因此研究重点不是单纯复现某一条现成方案，而是在动态剪枝、双向隐私与工程闭环之间完成兼容性重写。
